
\section{Results}
\subsection{Headline test-set performance}
All numbers read programmatically from committed JSONs at assembly time;
the official test partition of $n=624$ films was evaluated
exactly once per model. Table~\ref{tab:pne-main} gives all six
configurations.
\begin{table}[h]\centering\small
\begin{tabular}{lcccccc}
\hline
Model & Acc\% & Prec\% & Sens\% & Spec\% & F1\% & ROC-AUC \\
\hline
CNN core (baseline) & 72.12 & 69.35 & 99.23 & 26.92 & 81.65 & 0.9331 \\
RegionGCN (baseline) & 79.01 & 75.05 & 99.49 & 44.87 & 85.56 & 0.9514 \\
CNN core (class-weighted) & 81.73 & 79.11 & 96.15 & 57.69 & 86.81 & 0.9174 \\
RegionGCN (class-weighted) & 84.13 & 80.38 & 98.72 & 59.83 & 88.61 & 0.9471 \\
CNN core (census-cleaned) & 79.81 & 76.19 & 98.46 & 48.72 & 85.91 & 0.9278 \\
RegionGCN (census-cleaned) & 84.94 & 82.17 & 96.92 & 64.96 & 88.94 & 0.9377 \\
\hline
\end{tabular}
\caption{Pneumonia test-set metrics, all configurations (n=624 official held-out films).}\label{tab:pne-main}
\end{table}

\subsection{External head-to-head: TorchXRayVision on the same films}
The strongest external tool we could run unchanged on the identical test
set is TorchXRayVision's DenseNet-121 (\texttt{densenet121-res224-all},
published weights, official preprocessing). On the same 624 films it scores
accuracy 38.5\% and ROC-AUC 0.7990
(\texttt{results/pneumonia/test\_probs\_torchxrayvision.json}). Every one of our six configurations outperforms it on the same data and split --- the cleaned RegionGCN by +0.139 AUC and +46.5 threshold-accuracy points (threshold accuracy at 0.5 understates the transfer model's ranking ability, as its AUC shows, so we lead with AUC). We read this as a domain-shift finding --- adult-pretrained CXR models do not transfer zero-shot to this paediatric benchmark --- not as a general ``small beats big'' claim. The comparison stays like-for-like: identical films, identical split, tool used as published.
\subsection{The published-benchmark gap, stated honestly}
Kermany et al.\ report 92.8\% accuracy \cite{kermany2018} with a
transfer-learned Inception-v3 at 299px trained at substantially larger
effective scale; that number is paywalled and was corroborated through two
independent secondary sources (\texttt{results/kermany\_number\_verification.json}).
Our compact 2-CPU models do not close that cross-regime gap, and we do not
claim otherwise. The honest framing: on the strongest \emph{same-split} external tool comparison available, dataset-specific training prevails over zero-shot adult-domain transfer; against the
published cross-regime reference the gap is recorded as a redirected angle
(transfer initialisation from a large-scale pretrained trunk), pursued in
the transfer study below.
\subsection{Label-noise census}
The census (\texttt{results/pneumonia/label\_noise\_census.json}) is
admissible under the gate (OOF accuracy 89.81\%
vs.\ majority 74.20\%). Estimated noise rate
\textbf{21.15\%} over $n=4710$ training
films, 996 flagged identifiers, published
verbatim in the appendix. The asymmetry is again one-directional: 973
films labelled PNEUMONIA cross-validate as NORMAL, vs.\ 23 the other way.
\subsection{Cleaning and retraining}
Removing the 155 flagged films and retraining gives CNN
79.81\% and RegionGCN
\textbf{84.94\%} / AUC 0.9377
(\texttt{results/pneumonia/cleaned\_results.json}) --- the cleaned
RegionGCN is the strongest configuration in the study, a
+5.9-point gain over
its own baseline, and the class-weighted CNN gains
+9.6 points. Here
cleaning demonstrably helps, unlike the malaria regime; the difference is
discussed in the limitations section.
